\section*{章末练习}
\begingroup\small
\subsection*{口述概念}
\begin{enumerate}[nosep,leftmargin=*]
  \item 解释 $\sigma$ 代数、测度、可测函数、简单函数与 Lebesgue 积分的依赖顺序；为什么定义可测函数时使用逆像？
  \item 区分 Borel 可测集、Lebesgue 可测集、零测集与“几乎处处”。为什么 a.e. 结论不等于样本中永不失败？
  \item 比较 Fatou 引理、单调收敛定理和支配收敛定理的假设与结论，并解释共同支配函数为什么必须可积。
\end{enumerate}

\subsection*{笔试推导}
\begin{enumerate}[nosep,leftmargin=*]
  \item 从可数可加性证明测度的单调性、从下连续性；再证明从上连续性并指出 $\mu(A_1)<\infty$ 不能删除。
  \item 推导固定简单函数 $\phi_m$ 的积分公式 $1/2-2^{-(m+1)}$，并用单调收敛定理得到 $\int_0^1x\,dx=1/2$。
  \item 用 Fatou 引理分别作用于 $g+f_n$ 与 $g-f_n$，完成支配收敛定理的积分夹逼证明。
\end{enumerate}

\subsection*{数值编程}
\begin{enumerate}[leftmargin=*]
  \item 计算 $m=2,4,8$ 时逼近 $x$ 的二进台阶积分。
  \item 改变尖峰 $n\mathbf1_{(0,1/n]}$ 的 $n$，比较面积与固定点值。
  \item 比较两种 Pareto 模型的截断期望，并解释能否用 DCT。
\end{enumerate}

\subsection*{研究判断}
\begin{enumerate}[nosep,leftmargin=*]
  \item 研究员看到损失 $L_n(\omega)\to L(\omega)$ a.e.，便宣布期望收敛。列出可以补足结论的两条不同路线及各自证据。
  \item 一组收益数据只有有限个缺失值。能否直接把缺失集合称为零测集并忽略？分别从经验数据与概率模型层评价。
  \item 模拟程序在越来越密的网格上给出稳定积分。设计一份比“再加网格”更强的可积性与极限交换审计。
\end{enumerate}
\endgroup

\subsection*{基础衔接：独立计算}
设 $X$ 的密度为 $2x^{-3}$（$x\geq1$）。求 $\mathbb EX$ 与 $\mathbb EX^2$，判断能否用 $X^2$ 支配函数列并调用 DCT。
